% !TEX root = chapter-standalone.tex

\chaptersecond{Categorical ways of thinking, I}{chapter-rosti}{chap:diagrams}{
    In this chapter we discuss the religion that category theorists follow. 
}
% \devel{
%     The relation between graphs and categories might be confusing at first---there is a lot of overlap.
%     Both are composed of ``points and arrows''.
%     A graph is best thought of a \emph{data structure}: a list of nodes and edges.
%     This data structure could be used to describe a category.
%     But a category is something more: it has a closure property (all arrows compositions exist) and a constraint (composition is associative).
%     A graph needs not have these properties.
%     \todotext{\alphubel: It is weird to say it here because we haven't introduced categories.}
% }
\label{chap:diagrams}

\subimport{./}{10_up_to_isomorphism}
\subimport{./}{20_using_diagrams}
\subimport{./}{30_context}
\subimport{./}{40_commentary}
